% GENERATED FILE. Edit canonical lessons, then run npm run generate:books.
% canonical-source-hash: fnv1a64:583f6882da98811c
% canonical-lessons: LA-C366-videlicet, LA-C366-nimirum, LA-C366-atqui, LA-C366-proinde, LA-C366-idcirco

\chapter{Evidently, And Yet, So Then}
\label{ch:la-evidently-and-yet-so-then}
\hlchaptermodality{366}

\begin{chapteropening}
\textbf{By the end of this chapter:} \emph{I can say evidently and no doubt, object with and yet, and draw a consequence with so then and for that reason.}

Complete the last lesson of chapter 366: I can say evidently and no doubt, object with and yet, and draw a consequence with so then and for that reason.
\end{chapteropening}

\section[vidēlicet]{vidēlicet — evidently, clearly; namely}

\label{lesson:LA-C366-videlicet}



\begin{quote}
Before the new one: say the Latin for surely, then the Latin for to be sure.
\end{quote}

\subsection*{You'll want to know: vidēlicet}
\textbf{vidēlicet} — \textquotedblleft{}evidently, clearly; namely\textquotedblright{}.

\begin{quote}
Domum rediit, vidēlicet fessus.
\end{quote}

\begin{quote}
\emph{He went home, tired, evidently.}
\end{quote}

\emph{Vidēlicet} = \emph{vidēre licet}, one may see. English \emph{viz.} is its abbreviation.

\begin{cousinweb}
\emph{vidēlicet} is contracted from \emph{vidēre licet} \textquotedblleft{}one may see\textquotedblright{} (Lewis \& Short). English \emph{videlicet} and its short form \emph{viz.} were borrowed from it (EtymDB).
\end{cousinweb}

\begin{grammarlens}[title={What you've built}]
One more word for telling, explaining and describing at B1.
\end{grammarlens}

\subsection*{Your turn}
\begin{itemize}
\raggedright
  \item \emph{Say it:} \emph{vidēlicet}
  \item \emph{Say it:} \emph{vidēlicet}, once more
  \item \emph{Say it:} a sentence of your own with \emph{vidēlicet}
  \item \emph{Recall:} say the Latin for surely, then the Latin for to be sure, then say \emph{vidēlicet} again
\end{itemize}

\begin{tcolorbox}[breakable,colback=teal!4,colframe=teal!35!black,title={Before you move on}]
What does vidēlicet mean? (\textquotedblleft{}Evidently, clearly; namely\textquotedblright{}.) Say it once more.
\end{tcolorbox}
\section[nīmīrum]{nīmīrum — without doubt, surely}

\label{lesson:LA-C366-nimirum}



\begin{quote}
Before the new one: say the Latin for to be sure, then the Latin for evidently.
\end{quote}

\subsection*{You'll want to know: nīmīrum}
\textbf{nīmīrum} — \textquotedblleft{}without doubt, surely\textquotedblright{}.

\begin{quote}
Nīmīrum vērum dīcis.
\end{quote}

\begin{quote}
\emph{No doubt you are telling the truth.}
\end{quote}

\emph{Nīmīrum} = \emph{nī mīrum}, no wonder. It can carry a smile of irony.

\begin{cousinweb}
\emph{nīmīrum} is built from \emph{nī} \textquotedblleft{}not\textquotedblright{} and \emph{mīrum} \textquotedblleft{}a wonder\textquotedblright{} (EtymDB; Lewis \& Short): it is no wonder, so surely.
\end{cousinweb}

\begin{grammarlens}[title={What you've built}]
One more word for telling, explaining and describing at B1.
\end{grammarlens}

\subsection*{Your turn}
\begin{itemize}
\raggedright
  \item \emph{Say it:} \emph{nīmīrum}
  \item \emph{Say it:} \emph{nīmīrum}, once more
  \item \emph{Say it:} a sentence of your own with \emph{nīmīrum}
  \item \emph{Recall:} say the Latin for to be sure, then the Latin for evidently, then say \emph{nīmīrum} again
\end{itemize}

\begin{tcolorbox}[breakable,colback=teal!4,colframe=teal!35!black,title={Before you move on}]
What does nīmīrum mean? (\textquotedblleft{}Without doubt, surely\textquotedblright{}.) Say it once more.
\end{tcolorbox}
\section[atquī]{atquī — and yet, but (raising an objection)}

\label{lesson:LA-C366-atqui}



\begin{quote}
Before the new one: say the Latin for evidently, then the Latin for without doubt.
\end{quote}

\subsection*{You'll want to know: atquī}
\textbf{atquī} — \textquotedblleft{}and yet, but (raising an objection)\textquotedblright{}.

\begin{quote}
Dīcis tē pauperem esse; atquī vīllam magnam habēs.
\end{quote}

\begin{quote}
\emph{You say you are poor; and yet you have a large country house.}
\end{quote}

\emph{Atquī} brings in a fact that tells against what was said. It takes the indicative.

\begin{cousinweb}
\emph{atquī} is built from \emph{at} \textquotedblleft{}but\textquotedblright{} and \emph{quī}, an old ablative form meaning \textquotedblleft{}in any way\textquotedblright{} (Lewis \& Short): but in any case.
\end{cousinweb}

\begin{grammarlens}[title={What you've built}]
One more word for telling, explaining and describing at B1.
\end{grammarlens}

\subsection*{Your turn}
\begin{itemize}
\raggedright
  \item \emph{Say it:} \emph{atquī}
  \item \emph{Say it:} \emph{atquī}, once more
  \item \emph{Say it:} a sentence of your own with \emph{atquī}
  \item \emph{Recall:} say the Latin for evidently, then the Latin for without doubt, then say \emph{atquī} again
\end{itemize}

\begin{tcolorbox}[breakable,colback=teal!4,colframe=teal!35!black,title={Before you move on}]
What does atquī mean? (\textquotedblleft{}And yet, but (raising an objection)\textquotedblright{}.) Say it once more.
\end{tcolorbox}
\section[proinde]{proinde — therefore, accordingly; so then}

\label{lesson:LA-C366-proinde}



\begin{quote}
Before the new one: say the Latin for without doubt, then the Latin for and yet.
\end{quote}

\subsection*{You'll want to know: proinde}
\textbf{proinde} — \textquotedblleft{}therefore, accordingly; so then\textquotedblright{}.

\begin{quote}
Hostēs adsunt; proinde arma capiāmus.
\end{quote}

\begin{quote}
\emph{The enemy is here; so then let us take up arms.}
\end{quote}

\emph{Proinde} draws a practical consequence and often comes with a subjunctive of exhortation (\emph{capiāmus}, let us take) or with a command.

\begin{cousinweb}
\emph{proinde} is built from \emph{prō} \textquotedblleft{}for, in proportion to\textquotedblright{} and \emph{inde} \textquotedblleft{}from there\textquotedblright{} (Lewis \& Short): from that, accordingly.
\end{cousinweb}

\begin{grammarlens}[title={What you've built}]
One more word for telling, explaining and describing at B1.
\end{grammarlens}

\subsection*{Your turn}
\begin{itemize}
\raggedright
  \item \emph{Say it:} \emph{proinde}
  \item \emph{Say it:} \emph{proinde}, once more
  \item \emph{Say it:} a sentence of your own with \emph{proinde}
  \item \emph{Recall:} say the Latin for without doubt, then the Latin for and yet, then say \emph{proinde} again
\end{itemize}

\begin{tcolorbox}[breakable,colback=teal!4,colframe=teal!35!black,title={Before you move on}]
What does proinde mean? (\textquotedblleft{}Therefore, accordingly; so then\textquotedblright{}.) Say it once more.
\end{tcolorbox}
\section[idcircō]{idcircō — for that reason, therefore}

\label{lesson:LA-C366-idcirco}



\begin{quote}
Before the new one: say the Latin for and yet, then the Latin for therefore.
\end{quote}

\subsection*{You'll want to know: idcircō}
\textbf{idcircō} — \textquotedblleft{}for that reason, therefore\textquotedblright{}.

\begin{quote}
Aeger erat; idcircō nōn vēnit.
\end{quote}

\begin{quote}
\emph{He was ill; for that reason he did not come.}
\end{quote}

\emph{Idcircō} points ahead to a reason or an aim: \emph{idcircō … quod}, for this reason, that; \emph{idcircō … ut}, for this purpose, that.

\begin{cousinweb}
\emph{idcircō} is built from \emph{id} \textquotedblleft{}that\textquotedblright{} and an old form of \emph{circā}, \emph{circum} \textquotedblleft{}around, about\textquotedblright{} (Lewis \& Short give \emph{id-circa}, EtymDB \emph{id} + \emph{circum}): about that, for that reason.
\end{cousinweb}

\begin{grammarlens}[title={What you've built}]
One more word for telling, explaining and describing at B1.
\end{grammarlens}

\subsection*{Your turn}
\begin{itemize}
\raggedright
  \item \emph{Say it:} \emph{idcircō}
  \item \emph{Say it:} \emph{idcircō}, once more
  \item \emph{Say it:} a sentence of your own with \emph{idcircō}
  \item \emph{Recall:} say the Latin for and yet, then the Latin for therefore, then say \emph{idcircō} again
\end{itemize}

\begin{tcolorbox}[breakable,colback=teal!4,colframe=teal!35!black,title={Before you move on}]
What does idcircō mean? (\textquotedblleft{}For that reason, therefore\textquotedblright{}.) Say it once more.
\end{tcolorbox}
